
Retrieval-augmented generation systems commonly filter corpora using GPT-2 perplexity, a pretraining-derived quality signal optimizing for narrative fluency. Whether retrieval performance requires different quality signals remains untested. We trained classifiers on BEIR retrieval success examples using stratified sampling to enforce independence from educational quality, then evaluated whether the resulting corpora improved retrieval performance and whether improvements derived from increased factual entity density. In proof-of-concept validation on BEIR Natural Questions, retrieval-quality filtering achieved Recall@10 of 0.520 versus a perplexity baseline of 0.470 (delta +0.050, +10.6\% relative improvement, exceeding the +0.03 gate threshold). However, mechanism testing refuted the hypothesized causal pathway: retrieval-selected documents exhibited 2.7\% lower entity density than perplexity-matched baselines (ratio 0.973 vs. target $\geq 1.15)$, and showed no differential advantage on semantic queries requiring paraphrasing ($\Delta Recall$ = 0.00pp vs. target $\geq 4pp)$. The proof-of-concept validation demonstrates pipeline feasibility and suggests retrieval-quality signals may diverge from pretraining quality, while falsifying entity density as the operative mechanism. Confirmation with full corpus-scale retrieval on Common Crawl is recommended. Our negative mechanistic findings redirect future work toward alternative explanations including query-document semantic alignment, answer-bearing sentence structure, or non-entity informativeness.
